\documentclass[11pt,a4paper]{article}
\usepackage[utf8]{inputenc}
\usepackage[margin=1in]{geometry}
\usepackage{amsmath,amssymb,amsfonts}
\usepackage{hyperref}
\usepackage{listings}
\usepackage{xcolor}
\usepackage{booktabs}

% Custom listing configuration for Tamarin code
\definecolor{keywordcolor}{rgb}{0.5, 0.0, 0.5}
\definecolor{stringcolor}{rgb}{0.13, 0.54, 0.13}
\definecolor{commentcolor}{rgb}{0.5, 0.5, 0.5}
\definecolor{backcolor}{rgb}{0.97, 0.97, 0.97}

\lstdefinelanguage{Tamarin}{
  keywords={theory, begin, end, builtins, functions, equations, rule, lemma, all-traces, exists-trace, left, right, diff},
  keywordstyle=\color{keywordcolor}\bfseries,
  ndkeywords={In, Out, Fr, K},
  ndkeywordstyle=\color{blue}\bfseries,
  identifierstyle=\color{black},
  sensitive=true,
  comment=[l]{//},
  morecomment=[s]{/*}{*/},
  commentstyle=\color{commentcolor}\itshape,
  stringstyle=\color{stringcolor},
  morestring=[b]",
  morestring=[b]'
}

\lstset{
  language=Tamarin,
  backgroundcolor=\color{backcolor},
  basicstyle=\ttfamily\small,
  breaklines=true,
  captionpos=b,
  keepspaces=true,
  numbers=left,
  numberstyle=\tiny\color{commentcolor},
  showspaces=false,
  showstringspaces=false,
  showtabs=false,
  tabsize=2
}

\title{\textbf{Tamarin Prover Notes: Protocol Specification, Logic, and Auto-Sources}}
\author{Study Notes}
\date{\today}

\begin{document}

\maketitle

\section{Specification of Cryptographic Protocols}

Tamarin models protocols as \textbf{Labelled Multiset Rewriting Systems (LMSR)} operating over a term algebra equipped with equational theories.

\subsection{Terms and Term Algebra}
Messages in Tamarin are represented as terms over a signature consisting of function symbols and variables:
\begin{itemize}
    \item \textbf{Variables and Types}: Fresh variables (\texttt{\textasciitilde x}), Public variables (\texttt{\$x}), and untyped variables (\texttt{x}).
    \item \textbf{Built-in Equational Theories}: Activated using the \texttt{builtins} construct.
    \item \textbf{User-defined Equations}: Declared using \texttt{functions} and \texttt{equations} keywords.
\end{itemize}

\subsection{State and Facts}
System state is represented as a multiset of Facts:
\begin{itemize}
    \item \textbf{Linear Facts}: Consumed when matched in a rule premise (e.g., \texttt{StateA1(A, s)}).
    \item \textbf{Persistent Facts}: Prefixed with \texttt{!}. Retained permanently after being produced (e.g., \texttt{!Pk(A, pk(sk))}).
    \item \textbf{Built-in Special Facts}:
    \begin{itemize}
        \item \texttt{Fr(x)}: Fresh name generated by the execution environment.
        \item \texttt{In(m)}: Message received from the untrusted network.
        \item \texttt{Out(m)}: Message sent to the untrusted network.
    \end{itemize}
\end{itemize}

\subsection{Multiset Rewriting Rules}
Rules specify state transitions and have the general form:
\[
\text{rule Name:} \quad [l] \xrightarrow{\quad a \quad} [r]
\]
where $l$ is a multiset of premises, $a$ is a multiset of action labels (trace annotations), and $r$ is a multiset of conclusions.


\subsection{Adversary Model (Dolev-Yao)}
The network is fully controlled by the Dolev-Yao adversary:
\begin{itemize}
    \item The adversary intercepts all \texttt{Out(m)} facts, converting them into knowledge facts \texttt{K(m)}.
    \item The adversary can construct new messages from known terms using equational theories and feed them into any \texttt{In(m)} premise.
\end{itemize}


\section{Specification of Security Properties}

Properties are specified in a first-order logic evaluated over the trace of action labels produced during rule executions.

\subsection{Syntax}
Formulas combine action labels, temporal variables (prefixed with \texttt{\#}), term variables, and standard logical connectives:
\begin{itemize}
    \item \textbf{Action Predicate}: \texttt{Action(t1, t2) @ \#i} states that \texttt{Action(t1, t2)} occurred at timepoint \texttt{\#i}.
    \item \textbf{Temporal Order}: \texttt{\#i < \#j} or \texttt{\#i = \#j}.
    \item \textbf{Quantifiers}: \texttt{All} ($\forall$), \texttt{Ex} ($\exists$).
    \item \textbf{Connectives}: \texttt{\&} ($\land$), \texttt{|} ($\lor$), \texttt{==>} ($\implies$), \texttt{not} ($\neg$).
\end{itemize}

\subsection{Types of Security Lemmas}

\subsubsection{Executability / Sanity Checks}
Verified using the \texttt{exists-trace} attribute to ensure that the protocol can reach a desired state without contradictions.


\subsubsection{Secrecy, agreements etc}
Verified using the default \texttt{all-traces} quantification (holds for all valid execution traces).


\section{Advanced Features and Modeling Techniques}

\subsection{Oracles and Proof Heuristics}
When Tamarin's automated constraint solver encounters non-termination or deep search paths, custom heuristics can be supplied:
\begin{itemize}
    \item \textbf{Interactive Mode}: Step-by-step constraint resolution via web interface (\texttt{tamarin-prover interactive .}).
    \item \textbf{Oracle Scripts}: External Python scripts passed via \texttt{--oracle=script.py} that dynamically prioritize goal selection based on rank values.
\end{itemize}


\subsection{Sources Lemmas}

\subsubsection{Core Concept}
A \textbf{sources lemma} (or typing lemma) is a specialized helper theorem tagged with \texttt{[sources]}. It bounds backward search during proof exploration by establishing structural invariants on protocol messages.

\begin{itemize}
    \item \textbf{Purpose:} Resolves partial deconstructions and prevents non-termination (infinite loops) in Tamarin's constraint solver.
    \item \textbf{Execution Order:} Tamarin proves \texttt{[sources]} lemmas \emph{first} via induction, using them as verified premises to prune infeasible traces for subsequent security properties.
\end{itemize}

\subsubsection{Logical Structure}
A sources lemma typically asserts that whenever an honest rule consumes a tagged message or variable, that value must have either originated from a legitimate protocol step or been known to the adversary (\texttt{KU}).

\begin{equation*}
\forall m, i.\ \text{Recv}(m) @ i \implies \big(\exists j.\ \text{Created}(m) @ j \land j < i\big) \lor \big(\exists k.\ \text{KU}(m) @ k \land k < i\big)
\end{equation*}

\subsubsection{Practical Guidelines}
\begin{enumerate}
    \item \textbf{Auto-Sources:} Tamarin attempts automatic lemma generation via \texttt{auto-sources}. Manual definitions are required when complex nested structures cause non-termination.
    \item \textbf{Tagging Strategy:} Use explicit constant tags (e.g., \texttt{`1'}, \texttt{`3'}) in protocol tuples to prevent structural ambiguity and type flaw attacks.
    \item \textbf{Action Annotations:} Ensure key reception and creation events are properly annotated with action facts in your multiset rewrite rules.
\end{enumerate}

\subsection{Global Restrictions}

A \textbf{restriction} (\texttt{restriction}) is a global axiom applied to every trace generated by the protocol. If an execution path violates a restriction, Tamarin discards it immediately as invalid.

\begin{itemize}
    \item \textbf{Purpose:} Enforces unique nonces, strict execution phases, or mathematical axioms.
    \item \textbf{Limitations:} Unlike lemmas, restrictions are accepted as absolute truth without mathematical proof. Writing an incorrect restriction can accidentally rule out valid protocol behaviors or attacks, leading to false security guarantees.
\end{itemize}

\end{document}